% 将数据库的分类和排序字段交给 Biber，保留 ref.bib 原始数据。
\DeclareSourcemap{
  \maps[datatype=bibtex]{
    \map{\step[fieldsource=refsort,fieldtarget=sortkey]}
    \map{\pertype{newspaper}\step[fieldset=refgroup,fieldvalue={4}]}
    \map{\pertype{thesis}\pertype{mastersthesis}\pertype{phdthesis}
      \step[fieldset=refgroup,fieldvalue={5}]}
    \map{\pertype{online}\pertype{electronic}
      \step[fieldset=refgroup,fieldvalue={6}]}
    \map{\step[fieldsource=translator,final]\step[fieldset=refgroup,fieldvalue={9}]}
    \map{\pertype{book}\step[fieldset=refgroup,fieldvalue={7}]}
    \map{\step[fieldset=refgroup,fieldvalue={8}]}
  }
}
\DeclareSortingTemplate{report}{
  \sort[final]{\field{sortkey}}
  \sort{\field{sortname}\field{author}\field{editor}\field{sorttitle}\field{title}}
  \sort{\field{sortyear}\field{year}}
  \sort{\field{sorttitle}\field{title}}
  \sort{\field{entrykey}}
}
\ExecuteBibliographyOptions{sorting=report}

% 中文条目使用范例的紧凑标点。外文姓名与著录单元保留空格。
\newcommand*{\reportbibspace}{\iffieldequalstr{userd}{chinese}{}{\addspace}}
\newcommand*{\reportbibcomma}{\addcomma\reportbibspace}
\renewcommand*{\newunitpunct}{\addperiod\reportbibspace}
\renewcommand*{\subtitlepunct}{\addcolon\reportbibspace}
\renewcommand*{\publocpunct}{\addcolon\reportbibspace}
\renewcommand*{\pubdatadelim}{\setunit*{\reportbibcomma}}
\renewcommand*{\bibpagespunct}{\addcolon}
\DeclareDelimFormat[bib,biblist]{multinamedelim}{\reportbibcomma}
\DeclareDelimFormat[bib,biblist]{finalnamedelim}{\reportbibcomma}
\DeclareDelimFormat[bib,biblist]{andothersdelim}{%
  \iffieldequalstr{userd}{chinese}{}{\addcomma\addspace}}
\DeclareDelimFormat{titletypedelim}{\nobreak}
\expandafter\patchcmd\csname abx@macro@journal+issuetitle\endcsname
  {\setunit*{\addcomma\addspace}}{\setunit*{\reportbibcomma}}{}%
  {\PackageError{report-bibliography}{Cannot set journal punctuation}{}}
\expandafter\patchcmd\csname abx@macro@journal+issuetitle\endcsname
  {\setunit*{\addcomma\space}}{\setunit*{\reportbibcomma}}{}%
  {\PackageError{report-bibliography}{Cannot set volume punctuation}{}}

% DOI 保留在数据里。排版时也避免仅凭 DOI 追加 /OL。
\AtEveryBibitem{\clearfield{doi}}
\AtEveryCitekey{\clearfield{doi}}
\AtBeginBibliography{\interlinepenalty=10000\relax}

\forcsvlist{\DeclareBibliographyCategory}{ref1,ref2,ref3,ref4,ref5,ref6,ref7,ref8,ref9}
\AtDataInput{%
  \iffieldundef{refgroup}{}{%
    \csgdef{reportbibgroup@\thefield{refgroup}}{}%
    \addtocategory{ref\thefield{refgroup}}{\thefield{entrykey}}%
  }%
}

\makeatletter
% 脚注的完整著录复用文末的作者人数与姓名分隔符。
\DeclareCiteCommand{\report@fullcite}
  {\defcounter{maxnames}{\blx@maxbibnames}%
   \defcounter{minnames}{\blx@minbibnames}%
   \usebibmacro{prenote}}
  {\usedriver{\delimcontext{bib}\DeclareNameAlias{sortname}{default}}{\thefield{entrytype}}}
  {\multicitedelim}
  {\usebibmacro{postnote}}

\newcommand{\report@footcite}[1]{%
  \begingroup
    \setlength{\footnotesep}{8.5bp}%
    \def\thempfn{[\arabic{\@mpfn}]}%
    \def\@makefnmark{\textsuperscript{\@thefnmark}}%
    \footnote{\songti\fontsize{10.5pt}{10.5pt}\selectfont\report@fullcite{#1}.}%
  \endgroup
}
\renewrobustcmd*{\footcite}[1]{\forcsvlist{\report@footcite}{#1}}

\defbibenvironment{report}
  {\list{\mkbibbrackets{\printfield{labelnumber}}}{%
    \setlength{\leftmargin}{21bp}%
    \setlength{\labelwidth}{\leftmargin}%
    \setlength{\labelsep}{0pt}%
    \setlength{\itemsep}{0pt}%
    \setlength{\parsep}{0pt}%
    \setlength{\topsep}{0pt}%
    \setlength{\partopsep}{0pt}%
    \renewcommand*{\makelabel}[1]{\makebox[\labelwidth][l]{##1}}%
  }}
  {\endlist}
  {\item}

\defbibfilter{reportbooks}{category=ref1 or category=ref2 or category=ref9}
\defbibfilter{reportforeign}{category=ref7 or category=ref8}

\newcommand{\report@bibgroup}[2]{%
  \subsection{#1}%
  \printbibliography[heading=none,env=report,resetnumbers=true,#2]%
}
\newcommand{\PrintReportBibliography}{%
  \clearpage
  \begingroup
    \buaa@reportbibpage
    \pagestyle{postref}%
    \section{主要参考文献}%
    \thispagestyle{postref}%
    \renewcommand*{\bibfont}{\setstretch{1.5}\normalfont\songti\fontsize{10.5bp}{12bp}\selectfont}%
    \ifboolexpr{test {\ifcsdef{reportbibgroup@1}}
      or test {\ifcsdef{reportbibgroup@2}} or test {\ifcsdef{reportbibgroup@9}}}
      {\report@bibgroup{专著}{filter=reportbooks}}{}%
    \ifcsdef{reportbibgroup@3}{\report@bibgroup{期刊论文}{category=ref3}}{}%
    \ifcsdef{reportbibgroup@4}{\report@bibgroup{报纸文章}{category=ref4}}{}%
    \ifboolexpr{test {\ifcsdef{reportbibgroup@7}} or test {\ifcsdef{reportbibgroup@8}}}
      {\report@bibgroup{外文文献}{filter=reportforeign}}{}%
    \ifcsdef{reportbibgroup@5}{\report@bibgroup{学位论文}{category=ref5}}{}%
    \ifcsdef{reportbibgroup@6}{\report@bibgroup{电子资源}{category=ref6}}{}%
    \clearpage
  \endgroup
}
\makeatother
